\textbf{结论名称：}抛物线的定义（焦准定义）

\textbf{适用条件：}对于平面内一动点 \(P\)、一定点 \(F\)（焦点）和不过 \(F\) 的定直线 \(\ell\)（准线）；点 \(P\) 在抛物线上当且仅当定义成立。

\textbf{变量说明：}设焦点为 \(F\)，准线为 \(\ell\)，抛物线上点为 \(P\)，\(H\) 为 \(P\) 在 \(\ell\) 上的垂足；\(PF\) 表示点 \(P\) 到焦点 \(F\) 的线段长，\(PH\) 表示点 \(P\) 到直线 \(\ell\) 的距离。

\textbf{核心公式：}
\[
\boxed{PF = PH}
\]
或等价写作 \( \dfrac{PF}{\operatorname{dist}(P,\ell)} = 1 \)。

\textbf{使用场景：}
\begin{enumerate}
\item 已知焦点和准线求抛物线的标准方程；
\item 将焦半径长度转化为点到准线的距离，简化长度计算；
\item 用于证明与焦点弦、抛物线切线等相关的几何性质。
\end{enumerate}

\textbf{简单例子：}
已知抛物线焦点 \(F(2,0)\)，准线 \(\ell: x=-2\)。验证点 \(P(2,4)\) 在抛物线上：
\[
PF = \sqrt{(2-2)^2 + (4-0)^2} = 4,\qquad
PH = |\,2 - (-2)\,| = 4,
\]
满足 \(PF = PH\)。由定义可建立轨迹方程，整理得标准方程 \(y^2 = 8x\)。